\section{El Descenso a lo Circunstancial}\label{between-chapters-two-and-three}

Ambos capítulos sobre el tarot han funcionado en el marco de la estructura de 78 cartas: los Arcanos Mayores como fuerzas cosmológicas, los Arcanos Menores como acontecimientos sefirotico-elementales y las cartas de la corte como personas o formas de fuerza elemental. El registro ha sido, incluso en su vertiente más práctica, el de las fuerzas y los arquetipos.

El sistema de naipes inglés deja todo eso de lado. No hay Sefirot, ni letras hebreas, ni arquitectura cabalística. La pregunta a la que responde no es qué fuerza se mueve, sino qué está pasando: a quién, en qué ámbito, en qué dirección del tiempo. No son la misma pregunta, y el instrumento adecuado para una no es el adecuado para la otra.

Este es el sistema más fácil de usar del libro. Un lector con experiencia y una baraja estándar no necesita nada más para abordar la textura práctica y circunstancial de la vida de quien consulta. Su autoridad es totalmente pragmática: ha funcionado con suficiente consistencia, en suficientes manos y durante suficiente tiempo como para que se pueda confiar en él. Eso es suficiente.

\begin{center}\rule{0.5\linewidth}{0.5pt}\end{center}

\begin{center}\rule{0.5\linewidth}{0.5pt}\end{center}
